\chapter{Hardware}

\section{Die Schaltung des TransistorTesters}
\label{sec:hardware}
Die Schaltung des TransistorTesters in Abbildung~\ref{fig:ttester} basiert auf der Schaltung von
Markus F., die er in Abb. 1 des AVR-Transistortester Reports \cite{Frejek} veröffentlicht hat.
Geänderte oder verschobene Bauteile sind mit \textcolor{green}{grüner Farbe} markiert, optionale Teile sind
mit \textcolor{red}{roter Farbe} gekennzeichnet.

Einige Änderungen wurden gemacht, weil die Strom-Abschaltung in einigen Nachbauten Probleme
bereitet hatte.
Deshalb ist der Widerstand R7 auf \(3.3k\Omega\) reduziert. 
Der Kondensator C2 ist auf 10nF verkleinert und der Widerstand R8 ist verschoben, so dass der
Ausgang PD6 nicht versucht, den C2 Kondensator direkt zu laden.
Zusätzliche Abblock-Kondensatoren wurden hinzugefügt und sollten nahe den Versorgungs-Anschlüssen
des ATmega und nahe bei dem Spannungsregler plaziert werden.

Weil der PD7-Eingang und der PC6-Anschluss (RESET) die einzigen Anschlüsse sind, wo
,,pull-up'' Widerstände gebraucht werden, wurde ein zusätzlicher \(27k\Omega\) Widerstand am PD7 (Pin 13) vorgesehen.
Mit dieser Änderung können die internen ,,pull-up'' Widerstände des ATmega abgeschaltet werden.

Ein Quarz mit seinen 22pF Kondensatoren wurde zusätzlich vorgesehen.
Ein Quarz hat Vorteile für die Kapazitätsmessung wegen der genaueren Zeitmessung.

Die neue Software kann den Spannungsbereich für den ADC umschalten. Die Um\-schalt-Ge\-schwin\-dig\-keit
wird durch den externen Kondensator C1 am AREF-Pin (21) des ATmega reduziert.
Um die Messung nicht langsamer als notwendig machen zu müssen, sollte der Kondensator auf
1nF reduziert werden. Ein Entfernen des Kondensators ist ebenfalls möglich.
Zum Anpassen der Software an die jeweilige Schaltung schauen Sie bitte in dem
Kon\-fi\-gura\-tions-Kapitel~\ref{sec:config} nach. 

Einige unterschiedliche Kombinationen von R11 / R12 zirkulieren im Internet.
Ich habe die Software an den Original-Entwurf von Markus F.~\cite{Frejek} mit \(10k\Omega\) und \(3,3k\Omega\) angepasst.

Die zusätzliche 2,5V Präzisions-Spannungsreferenz, die an Pin PC4 (ADC4) angeschlossen ist,
wird für die Überprüfung und Kalibration der VCC Versorgungsspannung benutzt, ist aber nicht
erforderlich.
Sie können eine LM4040-AIZ2.5 (0.1\%),
eine LT1004CZ-2.5 (0.8\%) oder eine LM336-Z2.5 (0.8\%) Spannungsreferenz benutzen.
Ein zusätzlicher ISP-Anschluss wurde hinzugefügt, um leichter neue Software-Versionen
laden zu können.

\begin{figure}[H]
\centering
\includegraphics[width=18cm]{../FIG/ttester.eps}
\caption{Neue TransistorTester-Schaltung}
\label{fig:ttester}
\end{figure}

Um einen einfacheren Anschluss des Displays an den ATmega auf Streifenleiterplatinen zu erreichen,
kann die Software auch eine andere Port-D-Belegung berücksichtigen.
Die folgende Tabelle \ref{tab:grid-change} zeigt die Änderungen der Pinbelegungen.

\begin{table}[H]
  \begin{center}
    \begin{tabular}{| c | c | c |}
    \hline
       Signal & Schaltbild & Streifenleiter\\
    \hline
    Tastensignal  &  PD7   &  PD0 \\
    LCD-RS    &  PD4      & PD7 \\
    LCD-E     &  PD5   & PD5 \\
    LCD-D4    &  PD0   & PD4 \\
    LCD-D5    &  PD1   & PD3 \\
    LCD-D6    &  PD2   & PD2 \\
    LCD-D7    &  PD3   & PD1 \\
    \hline
    \end{tabular}
  \end{center}
  \caption{Änderungen für Streifenleiter-Platine}
  \label{tab:grid-change}
\end{table}

Zum besseren Schutz der ATmega-Eingänge kann eine Erweiterung mit einem Relais nach Schaltbild~\ref{fig:relay_addon}
angeschlossen werden. Die Ruhekontakte des Relais schützen den ATmega im spannungslosen Zustand.
Die Konkakte werden von der Software nur für die Messung freigegeben.

\begin{figure}[H]
\centering
\includegraphics[width=7cm]{../FIG/relay_addon.eps}
\caption{Zusätzlicher Schutz der ATmega Eingänge}
\label{fig:relay_addon}
\end{figure}

Wenn die Ausgabe der seriellen Texte nicht gebraucht wird, kann der Pin PC3 des ATmega für die Messung
einer externen Spannung benutzt werden. Die Spannung kann mit den optionalen 10:1 Widerstandsteiler
bis zu 50V betragen und kann für die Messung der Zener Spannung einer Diode benutzt werden.
Ein Strom begrenzendes Netzteil mit bis zu 50V Ausgangsspannung kann mit dem 0V Signal des ATmega PD7 Pins
eingeschaltet werden, um zum Beispiel die Durchbruchspannung von Zenerdioden zu messen.
Einen Vorschlag für diese Erweiterung zeigt Abbildung~\ref{fig:zener} .
Der Tester zeigt die externe Spannung so lange an, wie man den Taster gedrückt läßt.
Ungefähr 40mA mehr Batteriestrom wird für diese Erweiterung bei gedrückter Taste gebraucht.

\begin{figure}[H]
\centering
\includegraphics[width=12cm]{../FIG/zener_exp.eps}
\caption{Erweiterung zum Messen von Zener Spannungen}
\label{fig:zener}
\end{figure}

\section{Hinweise für den Aufbau des TransistorTesters}
Jede LCD-Anzeige-Einheit mit mindestens 2x16 Zeichen und einem HD44780-kompatiblen Kontroller kann mit
dem TransistorTester benutzt werden.
Man sollte auf den Strombedarf der Hintergrundbeleuchtung achten, einige Anzeigen benötigen
mehr Strom als andere.
Ich habe OLED Anzeigen ausprobiert, aber diese Anzeigen haben teilweise die Messung des
ATmega beeinflusst und werden nicht empfohlen. Auch das Laden der Spezialzeichen für die 
Widerstandsdarstellung hat mit OLED's Schwierigkeiten ergeben.

Die Widerstände R1 bis R6 sind kritisch für die Messungen und diese \(680\Omega\) und
\(470k\Omega\) Widerstände sollten Messwiderstände sein (Toleranz von 0,1\%), um 
die volle Genauigkeit zu erreichen.
Man sollte Präzisions-Sockel für den ATmega Mikrocontroller verwenden um
die Austauschbarkeit des Mikrocontrollers sicherzustellen.
Es kann ein ATmega8, ATmega168 und ATmega328 Mikrocontroller verwendet werden.
Empfohlen werden ATmega168 oder ATmega328, wenn man alle Funktionen nutzen möchte.

Jedenfalls sollte man zuerst alle Bauteile ohne den Mikrocontroller bestücken.
Es wird als IC2 ein moderner ,,low voltage drop'' Spannungsregler wie MCP1702-5002 empfohlen, weil
dieser nur \(2\mu A\) Ruhestrom benötigt und auch noch 5V liefern kann, 
 wenn die Eingangsspannung nur 5,4V beträgt.
Aber dieser Regler ist leider nicht Pin-kompatibel zum bekannteren 78L05 Regler im TO-92 Gehäuse!

Nachdem alle benötigten Bauteile bestückt sind, sollte zuerst die Batterie
oder das Netzteil angeschlossen werden. Das LCD sollte dabei nicht angeschlossen sein, und der Mikrocontroller sich noch nicht im Sockel befinden.
Man sollte die Betriebsspannung des Mikrocontrollers und der LCD-Anzeige
überprüfen während der Start-Taster gedrückt wird.
Die Betriebsspannung sollte verschwinden, wenn man den Start-Taster loslässt.
Wenn die Betriebsspannung die richtige Polarität und Grösse hatte,
sollte man die Strom-Versorgung entfernen und den Mikrocontroller 
richtig herum einstecken. Sind Sie bitte vorsichtig und stellen Sie sicher,
dass alle Pinne des Mikrocontrollers im Sockel stecken.
Danach können Sie das LCD anschliessen. Prüfen Sie, dass die GND- und VCC-Anschlüsse des LCD richtig mit der Baugruppe verbunden sind.

Wenn Sie sicher sind, dass alles richtig angeschlossen ist, schliessen Sie
die Spannungsversorgung wieder an.
Wenn Sie den ATmega schon programmiert haben, können Sie den Start-Taster
drücken.
Durch das Drücken des Start-Tasters sollte die Hintergrundbeleuchtung
der LCD-Anzeige angehen.
Wenn Sie den Taster loslassen, sollte die LED auf der Platine schwach leuchten.
Beachte, dass die Software für den Mikrocontroller für den richtigen
Prozessor-Typ übersetzt sein muss. Ein Programm für den ATmega8 läuft
nicht auf einem ATmega168!

\section{Umrüstung von Tester Versionen nach Markus F.}
\label{sec:change_markus}
\begin{description}

\item[Spannungsüberwachung]  
Das Problem zeigt sich durch sofortiges Abschalten beim Einschaltversuch.
Bei den von mir empfohlenen Einstellungen der fuses (Makefile) wird die Spannungsüberwachung der
verschiedenen ATmega Versionen auf 4V gesetzt (brown out level). Deswegen kann es beim
Einschalten des Testers zu Problemen kommen, weil der Pin PD6 versucht den 100nF Kondensator C2
direkt zu schalten. Dabei kann es zu einem unerwünschten 5V Spannungseinbruch kommen.
Der Kondensator C2 kann problemlos auf <10nF verkleinert werden. Nach Möglichkeit sollte
man statt der direkten Verbindung PD6 zum Kondensator einen Widerstand >\(220 \Omega\) 
als Verbindung benutzen.
\item[Verbessern des Einschaltverhaltens]
Der Fehler zeigt sich oft, daß der Tester bei gedrücktem Taster zwar einschaltet, aber wieder
abschaltet, wenn der Taster losgelassen wird. Das Problem tritt öfter auf, wenn die Hintergrund-
Beleuchtung des LCD viel Strom braucht.
Der Widerstand R7 zur Basis des PNP Transistors T3 war mit \(27k \Omega\) sehr auf Stromsparen
optimiert. Der Widerstand sollte besser auf \(3.3k \Omega\) verkleinert werden um auch bei
geringerer Batteriespannung oder bei geringem Stromverstärkungsfaktor des PNP Transistors T3
ein sicheres Einschalten zu gewährleisten.
\item[Zusätzlicher Pull-Up Widerstand an PD7]
Der Fehler zeigt sich dadurch, daß der Tester nach einer kurzen Anzeigezeit mit der Meldung
,,Timeout'' abschaltet. Die Software ist standardmäßig so konfiguriert (Option PULLUP\_DISABLE),
daß die internen Pull-Up Widerstände abgeschaltet sind.
Dadurch ist der Pegel am Pin PD7 nicht mehr definiert,
wenn er nicht durch den Taster oder T2 auf GND Potential geschaltet ist. Ein externer
Pull-Up Widerstand von \(27k \Omega\) nach VCC vermeidet diesen Fehler.
\item[Kondensator C1 am AREF Pin]
In vielen Entwürfen wird hier ein 100nF Kondensator verwendet, so auch im Entwurf vom Markus Frejek.
Solange die Referenzspannung des ADC nicht verändert wird, ist das auch in Ordnung.
Bei der Software für den Transistortester für den ATmega168/328 wird aber eine automatische
Umschaltung der Referenzspannung von 5V auf die interne Referenzspannung von 1.1V vorgenommen,
wenn die Eingangsspannung unter etwa 1V liegt. Damit wird eine bessere Auflösung erreicht.
Leider erfolgt die Umschaltung von 5V auf die 1.1V sehr langsam, was eine zusätzliche
Wartezeit von 10ms erfordert. Durch den Austausch des 100nF Kondensators gegen einen 1nF
kann die Wartezeit deutlich verringert werden. Einen Einfluß des kleineren Kondensators
auf die Qualität der Meßergebnisse habe ich nicht feststellen können. Selbst das Entfernen
des Kondensators hat keinen wesentlichen Einfluß. Wer den 100nF Kondensator unbedingt
beibehalten möchte, kann die Makefile Option NO\_AREF\_CAP entfernen, um die längere
Wartezeit im Programm zu aktivieren.
\item[Nachrüsten eines 8 MHz Quarz]
Mit etwas Geschick kann auf der Lötseite der Platine ein 8MHz Quarz direkt an PB6 und PB7
(Pin 9 und Pin 10) nachgerüstet werden.
Bei meiner Nachrüstung habe ich auf die beiden 22pF Kondensatoren verzichtet.
Bei allen eingesetzten Prozessoren hat diese Lösung problemlos funktioniert.
Aber die Nachrüstung ist nicht unbedingt erforderlich. Die Taktfrequenz sollte aber
wegen der besseren Zeitkonstanten-Auflösung (Kapazität) auf jeden Fall 8MHz betragen.
Die Frequenz 8Mhz ist auch bei RC-Oszillator Betrieb durch Setzen der fuses möglich.
\item[Abblocken der Betriebsspannung]
Im Original Schaltbild vom Markus F. ist nur ein 100nF Kondensator zum Abblocken der
VCC Spannung (5V) eingezeichnet. Das ist deutlich zu wenig. Es sollte sowohl ein
100nF in unmittelbarer Nähe des ATmega als auch ein 100nF in unmittelbarer Nähe
des Spannungsreglers vorhanden sein. Auch an den Eingang des Reglers gehört ein
100nF Kondensator. Zusätzliche \(10\mu F\) Kondensatoren (Elektrolyt oder Keramik) am Eingang und
Ausgang des Reglers können die Spannungsstabilität verbessern. Keramische \(10\mu F\)
Kondensatoren in SMD Bauform sind zum Nachrüsten meist besser geeignet und
haben üblicherweise einen niedrigeren ESR Wert.
\item[Auswahl des ATmega Prozessors]
Die Grundfunktion des Testers ist immer noch mit dem ATmega8 möglich.
Dabei ist der Programmspeicher nahezu zu 100\% benutzt.
Da die ATmega168 oder ATmega328 Prozessoren pinkompatibel zum ATmega8 sind,
kann der Austausch nur empfohlen werden. Mittlerweile sind die Preise für
den ATmega328 so günstig, daß eigentlich nichts mehr für den ATmega168 spricht.
Mit dem ATmega168/328 gewinnt man folgende Vorteile:
Selbsttestfunktion mit automatischem Abgleich.\\
Erhöhung der Meßgenauigkeit durch automatische Umschaltung der ADC Referenzspannung.\\
Messung von Induktivitäten, deren Widerstandswert < \(2100 \Omega\) ist.\\
Messung des ESR Wertes von Kondensatoren > \(0.18 \mu F\).\\
Die Auflösung der Widerstandsmessung unter \(10 \Omega\) beträgt \(0.01 \Omega\).\\
Der PC3 Pin kann für eine serielle Ausgabe genutzt werden.
\item[Fehlende Präzisionsreferenz]
Normalerweise sollte die fehlende Präzisionsreferenz auch bei unbeschaltetem PC4 Pin
erkannt werden. In diesem Fall wird keine VCC=x.xV Anzeige in Zeile 2 beim Einschalten
angezeigt. Falls es zu der Anzeige kommen sollte, hilft ein nach VCC geschalteter 
\(2.2k \Omega\) Widerstand am PC4 Eingang.

\section{Chinesische Nachbauten}
Der Tester wird in China nach meinem Kenntnisstand in zwei Versionen nachgebaut.
Die erste Variante ist der Nachbau des ersten Entwurfs von Markus F. ohne ISP Schnittstelle.
Der bestückte ATmega8 ist bei dieser Version gesockelt, kann also auch durch einen ATmega168/328 ausgetauscht werden.
Für diese Version gelten alle Hinweise des Unterkapitels \ref{sec:change_markus}.
Zusätzliche 100nF keramische Kondensatoren sollten in der Nähe des ATmega an die VCC-GND und
AVCC-GND Anschlüsse zur besseren Spannungsstabilisierung angebraucht werden.
Desweiteren muß beachtet werden, daß bei der Nachrüstung eines Quarzes der ISP-Programmer selbst
einen externen Takt zuführen muß.\\

Die zweite Variante ist weitgehend in SMD Technik aufgebaut. Auch der ATmega168 ist in einem 32TQFP
Gehäuse fest verbaut.
Dafür ist ein 10 poliger ISP Stecker für die Programmierung auf der Platine vorgesehen.
Ich habe die Version ,,2.1 2012/11/06'' analysiert. Ein Fehler ist die Bestückung des Bauteils ,,D1'',
welches eigentlich die 2.5V Präzisionsreferenz sein soll. Bestückt ist aber eine Zenerdiode.
Dieses Bauteil sollte entfernt werden. Hier kann eine Präzisionsreferenz wie LM4040AIZ2.5 oder
LT1004CZ-2.5 angeschlossen werden. Eine fehlende Präzisionsreferenz wird von der Software erkannt,
so daß sie nicht unbedingt erforderlich ist.
Bei meinem Exemplar war die Software Version 1.02k installiert. Der 10 polige ISP Stecker war nicht
bestückt und ich mußte zusätzlich eine Brücke von Pin 10 nach Pin 6 nachlöten. Mein Programmer hat eine GND
Verbindung an Pin 10 erwartet, der Tester hatte aber nur bei Pin 4 und Pin 6 eine GND Verbindung.
Die Beschriftung des ATmega168 war abgeschliffen und es gab keine Dokumentation zum Gerät.
Die Sicherheits Bits des ATmega waren gesetzt, so daß sich das Programm nicht auslesen ließ.
Ich konnte aber problemlos die Software Version 1.05k installieren.
Die gleiche Softwareversion 1.05k hat bei einem anderen Nutzer mit der China Version ,,2.2 2012/11/26'' Probleme
gemacht. Die Software 1.05k lief hier erst, als ein weiterer 100nF SMD Kondensator zwischen die Pins 18-AVCC
und 21-GND ergänzt wurde. Die Software 1.05k benutzt bei Wartezeiten den Schlafzustand des ATmega.
Deswegen wechselt der Strombedarf häufiger und der Spannungsregler wird mehr beansprucht.
Aufgefallen ist mir weiter, daß die VCC Spannung mit einem \(100nF\) keramischen Kondensator und mit
einem \(220\mu F\) elektrolytischen Kondensator in der Nähe des 78L05 Spannungsreglers abgeblockt ist.
Die 9V Spannungszufuhr ist mit den gleichen Kondensatoren abgeblockt, allerdings  am Emitter des PNP Transistor
(parallel zur Batterie), nicht direkt am Reglereingang.
Die Leiterbahnen vom ATmega zu den Testports sind teilweise sehr dünn, so daß ich einen Widerstand
von etwa \(100m \Omega\) pro Signalweg messen konnte. Dies ist wohl mit der Grund dafür, daß zwei
mit \(0 \Omega\) verbundene Pinne einen Widerstandswert von \(0.3 \Omega\) messen.
Bei der ESR-Messung kann dies normalerweise durch den Nullabgleich kompensiert werden.
Bei der Messung von Widerständen unter \(10 \Omega\) werden die im Selbsttest ermittelten Offsets 
ab der Softwareversion 1.07k berücksichtigt.

\end{description}

\section{Programmierung des Mikrocontrollers}
Ich gebe die Software für den Mikrocontroller in Quelltext heraus.
Die Entwicklung wurde mit dem Linux Betriebssystem (Ubuntu) gemacht
und wird gesteuert mit einer Makefile.
Die Makefile stellt sicher, dass die Software entsprechend der vorher in der Makefile 
eingestellten Optionen übersetzt wird. Schauen Sie bitte in die LiesMich.txt Datei
im Verzeichnis trunk/default und in das Konfigurations-Kapitel~\ref{sec:config}.
Das Ergebnis der Übersetzung hat die Dateierweiterung .hex und .eep.
Üblicherweise heißen die Dateien TransistorTester.hex und TransistorTester.eep .
Die .hex Datei enthält die Daten für den Programmspeicher (Flash) des ATmega-Prozessors.
Die .eep Datei enthält die Daten für den EEprom-Speicher des ATmega.
Beide Dateien müssen in den richtigen Speicher geladen werden.

Zusätzlich muss der ATmega mit den ,,fuses'' richtig konfiguriert werden.
Wenn Sie meine Makefile zusammen mit dem Programm avrdude \cite{avrdude} benutzen, brauchen Sie
keine genaue Kenntnis über die Einzelheiten der fuses.
Sie brauchen nur ,,make fuses'' aufrufen, wenn Sie keinen Quarz benutzen oder Sie
müssen ,,make fuses-crystal'' aufrufen, wenn Sie einen 8MHz Quarz auf der Baugruppe installiert haben.
Bei der ATmega168 Serie der Mikrocontroller können Sie alternativ auch
,,make fuses-crystal-lp'' aufrufen für den low power Quarz Betrieb.
Benutzen Sie niemals die Quarz-Variante, wenn Sie keinen 8MHz Quarz installiert haben.
Wenn Sie sich nicht sicher mit den fuses sind, lassen Sie diese erst einmal wie
vom Werk gesetzt und bringen Sie den Tester in diesem Zustand zum Laufen.
Es kann sein, dass das Programm zu langsam läuft, wenn Sie die für den 8MHz Betrieb 
erzeugten Programmdaten benutzen, aber das kann man später korrigieren!
Aber falsch gesetzte fuses können die spätere ISP-Programmierung verhindern.
Wenn Sie unter dem Windows Betriebssystem arbeiten, ist der leichteste Weg zum
richtig programmierten ATmega zu kommen, die Benutzung der WinAVR Paketes \cite{winavr1},\cite{winavr2}.
Mit meinem Patch \cite{winavr3} können Sie auch die Fuses mit der Makefile setzen.
Natürlich muß das avrdude Programm Ihren Programmer unterstützen und die Konfiguration muß in
der Makefile richtig angepaßt sein.

\section{Fehlersuche}
Bei den meisten Problemen werden Sie den Text auf dem LCD-Display vermissen.
Zuerst sollten Sie prüfen, ob die LED auf der Platine schwach leuchtet, wenn Sie den Start-Taster
loslassen.
\begin{description}

\item[Gerät schaltet nicht ein]  
Wenn die LED nicht leuchtet, aber die VCC-Spannung den richtigen Wert hat, wenn
man den Start-Taster gedückt hält, schaltet der Mikrocontroller nicht richtig ein.
Der Mikrocontroller sollte die Spannung behalten indem der Ausgang PD6 auf 5V
geschaltet wird, was üblicherweise als eine der ersten Aktionen getan wird.
Wenn man die Start-Taste gedrückt hält, bleibt die Spannung ohnehin eingeschaltet.
So können Sie mit gedrücktem Taster den Wert der VCC-Spannung und zusätzlich den Spannungswert am
Ausgang PD6 prüfen.
Wenn die VCC-Spannung den richtigen Wert (5V) hat, aber die Spannung am Ausgang
PD6 unter 4V ist, startet der Mikrocontroller nicht richtig.
Für diesen Fall sollten Sie prüfen, ob die Programmdaten für den Flashspeicher
für den richtigen Prozessor-Typ ist und ob der Prozessor richtig konfiguriert ist (fuses).
Wenn der ATmega den Ausgang PD6 auf 5V schaltet und die Betriebsspannung 
trotzdem nicht eingeschaltet bleibt, wenn man den Start-Taster loslässt, ist der
Grund schwieriger zu finden.
Zuerst kann man die LED kurzschliessen und es noch einmal versuchen.
Wenn der Tester jetzt startet, ist die LED möglicherweise falsch herum eingebaut.
Wenn das nicht die Ursache ist, könnte der Grund ein unzureichender Stromverstärkungs-Faktor
des Transistors T3 (BC557C) sein.
Der Strom in die Basis von T3 ist niedriger, wenn der Mikrocontroller mit der LED einschaltet
wie im ,,Taster gedrückt''-Zustand.

\item[Nichts ist lesbar auf der LCD-Anzeige] 
Prüfen Sie die Spannung am Kontrast-Pin der LCD-Anzeige (Pin 3).
Stellen Sie mit dem Trimmer den Wert auf einen im Datenblatt angegebenen Wert und optimieren Sie
durch Sichtkontrolle.
Wenn Sie ein Hochtemperatur-Display haben, brauchen Sie eine negative Kontrast-Spannung für
den Betrieb.
In diesem Fall kann man den ICL~7660 Baustein zum Erzeugen der negativen Spannung aus der
positiven 5V verwenden.

Wenn keine Anzeige auf dem LCD erkannt wird und wenn die Hintergrundbeleuchtung an ist,
sollten Sie die Spannungsversorgung trennen und alle vier Datenverbindungen sowie die 
beiden Steuersignale überprüfen.
Wenn alle Verbindungen in Ordnung sind, sehe ich als Ursache nur noch die Möglichkeit einer
falschen Zeitabfolge der Steuersignale.
Die Ursache hierfür kann sein, dass der LCD-Controller langsamer ist als es die Software
des ATmega erwartet. Es könnte auch sein, dass der ATmega auf der falschen Taktrate läuft.
Bitte überprüfen Sie für welche Taktrate die Software übersetzt ist und ob
die fuses des ATmega für diese Geschwindigkeit richtig gesetzt sind.
Sie finden die eingestellte Taktrate in der betreffenden Makefile.
Wenn der Tester ohne die Abschalt-Elektronik aufgebaut ist, kann man mit einer
an die Test Pins angeschlossenen LED testen, ob das Programm arbeitet.
Wenn die LED flackert, läuft das Programm. Der Fehler muß in diesen Fall an
dem Anschluß des LCD's liegen. 
\item[Einiges, aber nicht alles ist auf der LCD-Anzeige lesbar] 
Überprüfen Sie ob die .eep Daten in den EEprom Speicher des ATmega geladen wurden.
Wenn alle Programmdaten richtig geladen wurden, sollten Sie die Taktrate ihrer
Programmdaten (Makefile) und die ATmega-Einstellungen prüfen (fuses).

\item[Messung ist zu langsam und Kapazitäten werden um Faktor 8 zu klein gemessen.] 
Sie betreiben die Software, die für 8MHz übersetzt wurde mit einer Taktrate von 1MHz.
Bitte konfigurieren Sie den ATmega mit den fuses richtig.

\item[Die Messung ergibt seltsame Ergebnisse]  
Überprüfe ob der ISP-Programmierstecker noch verbunden ist.
Der ISP-Stecker sollte nicht während einer Messung eingesteckt bleiben.
Sehr oft ist der Grund für falsche Messergebnisse, dass die Software mit der
 AUTOSCALE\_ADC Option und der NO\_REF\_CAP Option übersetzt wurde, aber der
Kondensator am AREF-Pin hat immer noch einen Wert von 100nF.
Falsche Bestückung von Bauteilen können auch eine Ursache für Messfehler sein 
oder zurückgebliebene Flussmittelreste können die Messung stören.
Bitte prüfen Sie nach Möglichkeit mit der Selbsttest-Funktion der
TransistorTester Software.
Zu den Einzelheiten schauen Sie in das Selbsttest Kapitel \ref{sec:selftest}.

Anderenfalls prüfen sie Ihre Platine visuell und prüfen Sie die Widerstandswerte
mit einem Ohmmeter. Sie können die Pins des ATmega für diese Prüfung benutzen,
zum Beispiel können Sie der Widerstand R1 zwischen Pin 23 und Pin 14 messen.
Schauen Sie in das Schaltbild \ref{fig:ttester} für die Einzelheiten.
Man braucht den Mikrocontroller nicht zu entfernen, nur die Stromversorgung sollte
vorher getrennt werden.

\item[Der Tester schaltet nach 2 Sekunden Anzeigezeit aus]  
Dies ist dann der Fall, wenn der Pull-Up Widerstand am PD7 Eingang
fehlt oder der Taster dauernd gedrückt wird. 
Die Software schaltet die internen Pull-Up Widerstände ab, um eine Beeinflussung
der Meßergebnisse auszuschließen. Deswegen ist ein externer Widerstand erforderlich.

\end{description}
